\documentclass{beamer}
\usetheme{Madrid}
\usepackage[utf8]{inputenc}
\usepackage[spanish]{babel}

\title{PageRank para Grafos Dirigidos Generales}
\subtitle{Proyecto Final - Analisis de Algoritmos}
\author{Nombre del Estudiante}
\institute{Yachay Tech University}
\date{2026}

\begin{document}

\begin{frame}
  \titlepage
\end{frame}

\begin{frame}{Motivacion}
  \begin{itemize}
    \item PageRank mide importancia en grafos dirigidos.
    \item La primera version funcionaba sobre un grafo pequeno.
    \item La version final procesa grafos generales desde archivos CSV.
    \item La implementacion principal esta escrita en C++.
  \end{itemize}
\end{frame}

\begin{frame}{Definicion formal}
  Para un grafo dirigido $G=(V,E)$:
  \[
    r(v)=\frac{1-d}{n}+d\sum_{u\in In(v)}\frac{r(u)}{out(u)}+d\frac{D}{n}
  \]
  \begin{itemize}
    \item $d$: factor de amortiguamiento.
    \item $D$: masa de nodos colgantes.
    \item $r$: distribucion estacionaria.
  \end{itemize}
\end{frame}

\begin{frame}{Paradigma}
  \begin{itemize}
    \item Iteracion de potencias.
    \item Modelo de caminante aleatorio.
    \item Teleportacion para garantizar convergencia.
    \item Criterio de parada con norma $L_1$.
  \end{itemize}
\end{frame}

\begin{frame}{Implementacion C++}
  \begin{itemize}
    \item Entrada: lista de aristas \texttt{source,target}.
    \item Mapeo de etiquetas de nodos a enteros.
    \item Listas de predecesores para evitar matriz densa.
    \item Salida: ranking, score, in-degree y out-degree.
  \end{itemize}
\end{frame}

\begin{frame}{Complejidad}
  \begin{itemize}
    \item Por iteracion: $O(n+m)$.
    \item Total: $O(k(n+m))$.
    \item Memoria: $O(n+m)$.
    \item Peor caso controlado por \texttt{max\_iter}.
  \end{itemize}
\end{frame}

\begin{frame}{Analisis empirico}
  \begin{itemize}
    \item Dataset manual para inspeccion.
    \item Datasets sinteticos de 100 a 5000 nodos.
    \item Se mide tiempo, iteraciones y error final.
    \item Todos los datasets convergieron con $\varepsilon=10^{-10}$.
  \end{itemize}
\end{frame}

\begin{frame}{Resultados}
  \begin{itemize}
    \item 100 nodos, 600 aristas: 22 iteraciones.
    \item 1000 nodos, 6000 aristas: 24 iteraciones.
    \item 5000 nodos, 30000 aristas: 24 iteraciones.
    \item Con $d=0.50$: 16 iteraciones; con $d=0.95$: 27 iteraciones.
  \end{itemize}
\end{frame}

\begin{frame}{Conclusiones}
  \begin{itemize}
    \item PageRank combina teoria de grafos y cadenas de Markov.
    \item La implementacion final ya no depende de un grafo pequeno.
    \item El analisis teorico coincide con la estructura del codigo.
    \item El proyecto queda listo para extenderse con datasets reales.
  \end{itemize}
\end{frame}

\begin{frame}{Preguntas}
  \centering
  Gracias.
\end{frame}

\end{document}
